\documentclass[a4paper,titlepage,onecolumn,oneside,12pt,nenglish,titlepage=firstiscover]{scrbook}
% Um während der Entwurfsphase weniger leere Seiten zu haben, kann man oben "twoside" durch "oneside" ersetzen.

\usepackage{todonotes} % disable this package in final thesis!

\usepackage[english]{babel} % language
\usepackage[utf8]{inputenc} % encoding
\usepackage{amsmath,amsfonts,amssymb,amsthm} % mathematical expression, symbols, fonts, and the whole theorem environment engine
\usepackage{xcolor,graphicx} % graphics and colorings
\usepackage{tabularx,booktabs,multirow} % all about tables: booktabs gives the standard formatting, multirow allows cells to span over multiple rows.
%\usepackage{scrextend} % super powerful tool that helps a bit with everything
\usepackage[margin=0pt]{subcaption} % specific form for captions
\usepackage[autostyle]{csquotes} % helper for quoting

\usepackage{rotating} % if you want to rotate any object

\usepackage[backend=bibtex,
  style=alphabetic, % alphabetic ordered list
	hyperref=true, 
	natbib=true, % makes use of natbib
	backref=true, % additional writes for each bib-entry, on which pages its cites
	backrefstyle=three, % this style replaces e.g. 11,12,13 by 11-13
	sortcites=true, 
	maxbibnames=50,
	maxcitenames=3, % cites at most three authors---if more, only first author followed by et al. is written
 	firstinits=true, % nur Initialen anzeigen und auch nur nach diesen sortieren
% 	refsection=part,
% 	refsegment=chapter,
	isbn=false,url=true,doi=false] 
	{biblatex} % for references
	
% % if you want to change pageback references are stated, you can use the following:
% \DefineBibliographyStrings{english}{%
%   backrefpage = {see p\adddot},%
%   backrefpages = {see pp\adddot}%
% }
\setlength\parindent{0pt}

\bibliography{bibfile} % add the .bib file you use here

% please fill in the following data:
\newcommand{\mytitle}{Automatic Optimization of Level Parameters for RSHash}
\newcommand{\mysubtitle}{RSHash is a minimizer based fast, compressed hash table for storing DNA k-mers. It was recently introduced by the Reinert Group.}
\newcommand{\myname}{Bjarne David Jehn}
\newcommand{\thesisType}{Bachelor's Thesis}

\usepackage[pdfdisplaydoctitle,menucolor=orange!40!black,filecolor=magenta!40!black,urlcolor=blue!40!black,linkcolor=red!40!black,citecolor=green!40!black,colorlinks]{hyperref}
\hypersetup{pdftitle={\mytitle}, pdfauthor={\myname}}


%%=========== Timestamp on each page. Helpful to keep track of versions. 
%%=========== Remove in final thesis!!!
%\usepackage{scalefnt}
%\usepackage{scrtime}
%\usepackage{eso-pic}
%\AddToShipoutPicture{
%	\put(20,50){\textcolor[gray]{0.5}{\begin{tabular}{l} \thesisType{} \hspace{1cm} \myname \\ \mytitle{} (Compile date: \today/\thistime)\end{tabular}}}
%}
%%=========== End Timestamp


%=========== graphics, in particular tikz
\usepackage{tikz}
\usetikzlibrary{arrows,decorations.pathmorphing,decorations.pathreplacing,backgrounds,positioning,fit,matrix}
\usetikzlibrary{shapes,calc,arrows.meta,positioning}

%=========== tables from csv files (helpful for experiments)
\usepackage{pgfplots}
\usepackage{siunitx}
\usepackage{pgfplotstable}
\usepgfplotslibrary{groupplots}
\pgfplotsset{compat=1.18}
\usepackage{csvsimple}

%=========== Pseudocode
\usepackage{algorithm}
\usepackage{algpseudocode}
\algnewcommand\algorithmicinput{\textbf{Input:}}
\algnewcommand\algorithmicoutput{\textbf{Output:}}
\algnewcommand\Input{\item[\algorithmicinput]}
\algnewcommand\Output{\item[\algorithmicoutput]}
\floatname{algorithm}{Algorithm}

%=========== \cref is a powerfull command that automatically states the current environment (e.g. Theorem, Lemma, etc.)
\usepackage[sort&compress,nameinlink,noabbrev,capitalize]{cleveref}
\crefname{algorithm}{algorithm}{algorithms}
\Crefname{algorithm}{Algorithm}{Algorithms}


%\usepgfplotslibrary{external}
%\tikzexternalize

%===========theorem environment
\theoremstyle{plain} % typical theorem-style (bold header, italic font)
\newtheorem{theorem}{Theorem}[chapter]
\newtheorem{lemma}[theorem]{Lemma}
\newtheorem{corollary}[theorem]{Korollar}
\newtheorem{observation}[theorem]{Beobachtung}
\newtheorem{proposition}[theorem]{Proposition}
\newtheorem{rrule}{Reduktionsregel}[section]

\crefname{rrule}{Reduktionsregel}{Reduktionsregeln} % how cleverref (\cref) has to handle these environments
\Crefname{rrule}{RR}{RRs} % with \Cref us can use what you defined here---somtimes useful if you want to abbreviate environment when citing (e.g. in tables or figures)

\theoremstyle{definition} % typical definition-style (bold header, normal font)
\newtheorem{definition}[theorem]{Definition}

\theoremstyle{remark} % typical theorem-style (italic header, normal font)
\newtheorem{example}{Example}
\newtheorem*{remark}{Remark}
\newtheorem{reduction}{Reduction}
\theoremstyle{plain}

%===========problem definition evironment: gets three arguments (1) problem name #1 (2) input specification #2 (3) question specification #3
\newcommand{\problemdef}[3]{
  \begin{center}
    \begin{minipage}{0.95\textwidth}
      \noindent
      \textsc{#1}
      
      \vspace{2pt}
      \setlength{\tabcolsep}{3pt}
      \begin{tabularx}{\textwidth}{@{}lX@{}}
        \textbf{Eingabe:} 		& #2 \\
        \textbf{Frage:} 	& #3
      \end{tabularx}
    \end{minipage}
  \end{center}
}



%=========== Nützliche Abkürzungen
\newcommand{\NN}{\mathbb{N}} % natürliche Zahlen
\newcommand{\RR}{\mathbb{R}} % reelle Zahlen
\newcommand{\QQ}{\mathbb{Q}} % rationale Zahlen
\newcommand{\ZZ}{\mathbb{Z}} % ganze Zahlen

\newcommand{\bigO}{\mathcal{O}} % Landausymbol

\newcommand\abs[1]{\left|#1\right|} % Betragsfunktion
\newcommand\norm[2]{\left\lVert#1\right\rVert_{#2}} % norm
\newcommand\enorm[1]{\norm{#1}{2}} % euclidean norm

%=========== commands for minimizer distribution
\newcommand{\sequence}{\mathcal{S}}
\newcommand{\listOfSuperKmers}{\textbf{skmers}}
\newcommand{\maxThreshold}{T}

%=========== commands for MINLP
\newcommand{\minimizerLength}{m}
\newcommand{\threshold}{t}
\newcommand{\level}{l}
\newcommand{\minimizer}{x}
\newcommand{\kmer}{q}

\newcommand{\access}{acc}
\newcommand{\contains}{rank}
\newcommand{\rank}{rank}
\newcommand{\select}{sel}
\newcommand{\compare}{comp}

\newcommand{\timeForOperationOnStructure}[2]{T_{#1_{#2}}} %#1:operation - e.g. rank, #2:structure - e.g. R

\newcommand{\varAverageFrequency}{\textbf{f}_{\textbf{avg}}}
\newcommand{\averageFrequency}{f_{avg}}

\newcommand{\possibleMinimizerLengths}{\mathcal{M}}
\newcommand{\possibleThresholds}{\mathcal{T}}

\newcommand{\indicatorm}{\mu}
\newcommand{\indicatort}{\tau}

\newcommand{\varMinimizerLength}{\textbf{\minimizerLength}}
\newcommand{\varThreshold}{\textbf{\threshold}}

%=========== commands for subsampling
\newcommand{\numberOfSamples}{n}
\newcommand{\lengthOfSamples}{l}

\newcommand{\genomeText}{T}
\newcommand{\sampledText}{T^S}

%========== Title Formating
\titlehead{
\begin{minipage}{0.8\textwidth}
\textbf{Technische Universität Berlin}\\
Faculty IV: Electrical Engineering and Computer Science\\
\textbf{Freie Universität Berlin}\\
Department of Mathematics and Computer Science\\
Algorithmic Bioinformatics\\
\todo[inline]{welche uni angeben, welche reihenfolge..?}
\end{minipage}
\hfill
\begin{minipage}{0.16\textwidth}
\includegraphics[width=\textwidth]{./figures/tu-logo}
\end{minipage}
\vspace{3.5cm}
}

\title{\mytitle}
\subtitle{(\mysubtitle)\vspace*{3cm}}

\author{
%	\Large{\textbf{\thesisType{}}}\vspace{0.1cm}\\
	\textbf{\large{\myname}}\vspace{0.1cm}\\
	\large{Thesis submitted in fulfillment of the requirements for the degree}\\
	\large{\grqq Bachelor of Science\grqq (B. Sc.) in the field of Computer Science}
	\vspace*{3cm}\\
		\small{
			\begin{tabular}{rl}
			    First Reviewer:& Prof. Dr. Knut Reinert\\
			    \hspace{2cm}Second Reviewer:& Prof. Dr. Stephan Kreutzer\\
			    Supervisor:& Jonas Schulte-Mattler
			\end{tabular}
		}
	\vspace*{1cm}
}

\date{} % only month and year
\begin{document}

\selectlanguage{english}

\maketitle

\section*{Declaration of Authorship}
\noindent
I hereby declare that the thesis submitted is my own, unaided work, completed without any external help. Only
the sources and resources listed were used. All passages taken from the sources and aids used, either
unchanged or paraphrased, have been marked as such.\\

Where generative AI tools were used, I have indicated the product name, manufacturer, the software version
used, as well as the respective purpose (e.g. checking and improving language in the texts, systematic
research). I am fully responsible for the selection, adoption, and all results of the AI-generated output I use.\\

I have taken note of the Principles for Ensuring Good Research Practice at TU Berlin dated 15 February 2023.
\url{https://www.static.tu.berlin/fileadmin/www/10002457/K3-AMBl/Amtsblatt\_2023/Amtliches\_Mitteilungsblatt\_Nr.\_16\_vom\_30.05.2023.pdf}\\

I further declare that I have not submitted the thesis in the same or similar form to any other examination
authority.\\

\begin{tabular}[h]{lll} Berlin, & \underline{\hspace{5cm}} & \underline{\hspace{7cm}}\\
 & \small{Date} & \small{Signature}
\end{tabular}
\normalsize

\newpage
\section*{Abstract}
\todo[inline]{Draft}
RSHash is a static, compressed data structure for k-mers that answers lookup queries using a leveled minimizer scheme with thresholds to cut off frequent minimizers. Building it requires specifying a minimizer lengths and a threshold for every level used. Choosing optimal parameters to build RSHash for a given genetic sequence and specification of optimality remains an open problem as both its space consumption and the lookup time depend on these parameters because they affect the distribution of minimizers. To address this, we first define properties of these distributions which are specifically relevant to the problem and propose a mixed integer non linear program (MINLP) modeling the space consumption and lookup time of RSHash aiming to combine it with a solver to find optimal parameters. We evaluate the proposed MINLP for a selection of benchmark datasets and can show that the model is highly accurate regarding RSHashs space consumption but less so for the lookup time. Furthermore, time complexity to compute the distribution stats required to run the model grow exponentially with the number of considered levels which means a similar complexity to that of a brutefore approach. 
\section*{Zusammenfassung}
\todo[inline]{Passt so? => Ins Deutsche übersezten}


\newpage
\tableofcontents
\newpage

\chapter{Introduction} \label{chap:introduction}
\textbf{Genetic strands} are chains of the nucleobases adenine (A), cytosine (C), guanine (G), thymine (T) and uracil (U). Adenine, cytosine and guanine are present in both DNA and RNA whilst thymine only appears in DNA and uracil is specific to RNA. A genetic sequence can thus be represented as a string over an alphabet of length four meaning a compact representation of two bits needed to store each base. Despite this, genetic databases are often very large as sequences usually consist of several billion bases e.g. approximately $3.1 \cdot 10^9$ bases for the human genomes sequences. \\
The rapid growth of available genetic data due to better and faster sequencing techniques has led to a need for efficient algorithms for processing genetic sequences.\\

\textbf{K-mers} have emerged as an essential component of many approaches to encounter this issue. The k-mers of a sequence are all its substrings of length $k$, thus a sequence $S$ with $|S| \geq k$ contains exactly $|S|-k+1$ k-mers; for example, the sequence $CGACAT$ consists of the 3-mers $CGA,GAC,ACA,CAT$.\\
Uses for k-mers include k-mer counting, frequency analysis, sequence indexing, genome assembly, sequence comparison, taxonomic classification, phylogeny reconstruction, protein sequence searching and alignment as well the construction of k-mer databases \cite{moeckel_survey_2024}. 

\textbf{Sequence comparison} is used to compare a sequence against large genomic reference datasets; for example, to classify to what species a collected DNA sample belongs or to analyze whether a patients DNA contains indicators linked to certain diseases.\\
One approach to this is to check whether each k-mer in the sequence of interest is present in the reference dataset, also referred to as a lookup of the queried k-mers. On the other hand, a so called locate query answers not only whether queried k-mers are present but also returns all positions of the queried k-mers in the sequences. These positions are referred to as their offsets.\\

\textbf{A trade-off between space consumption and lookup time} is an inherent problem for data structures designed to answer lookup or locate queries. This can be demonstrated by considering the two extremes for data structures enabling a lookup or locate of a queried k-mer $q$ in a sequence $S$ with $|S| \gg k$ that doesn't contain any duplicate k-mers\footnote{For sequences containing duplicate k-mers, the lookup case does not change much. However, in the locate case it would then be required to find all offsets of $q$. For the compact example this would require to always scan through the whole sequence. For the fast example this would mean that the hash table entry for each k-mer would store a list of all its offsets.}:\\
The most space efficient representation of all k-mers in $S$ is the sequence itself needing approximately two bits per k-mer (bpk). A lookup/locate query of a $q$ would require scanning through the whole sequence, comparing every k-mer with $q$ and returning whether $q$ appears in the sequence (and its offset for the locate case). Both queries would terminate upon either finding a match or reaching the sequences end, requiring $O(|S|)$ time.\\
A very fast lookup/locate on the other hand could be achieved by storing the appearance of every k-mer in $S$ in a hash set (locate: hash table) enabling both queries in $O(1)$ time. For this approach every k-mer would have to be stored in the hash set/table leading to a space consumption of $2k$ bits per k-mer in the lookup case and  $2k+log_2(|S|)$ bits per k-mer to enable locate queries which require the offsets to be stored as well.\\
Since both a lookup time in $O(|S|)$ as well as $2k$ bits per k-mer are not viable for large genomes, there has been much work on data structures enabling queries in much less than $O(|S|)$ whilst also maintaining a certain compression.\\

\textbf{Using so called minimizers} is a common approach to achieve this. A minimizer of length $m$ is the smallest m-mer of a k-mer regarding a total order over all strings of length $m$. Using the lexicographical order for example to determine the minimizer of length $2$ in the $6$-mer $CGACAT$ results in $AC$ because regarding that order, this is the smallest $2$-mer in $CGACAT$.\\
Minimizers have several properties which render them highly interesting for this domain. Firstly they are considered to be sparse because consecutive k-mers are very likely to share the same minimizer. \Citet{zheng_improved_2020} show that in a random string $S$ with $|S| \gg k > m$ each minimizer of length $m$ is expected to cover an average of $(k-m+2)/2$ k-mers. These consecutive k-mers sharing the same minimizer are also referred to as super-k-mers. Even in a sequence only containing unique k-mers, several super-k-mers can have the same minimizer as $m < k$.\\
A second interesting property is the skewness of the distribution of minimizers in regards to their \textit{frequencies} which is the number of super-k-mers which have a specific minimizer. In a given sequence only a few  minimizers have high frequencies whilst most have small frequencies or do not appear at all. Some intuition for this property can be gained by considering an m-mer consisting only of $A$s versus one of $T$s. Using the lexicographical order the first will be the minimizer of every super-k-mer in which it appears whilst the second would only become the minimizer of a super-k-mer containing only $T$s.\\

\chapter{Previous Work}
Both the sparsity and skewness of minimizers have been observed by several previous works \cite{schulte-mattler_rshash_2026, pibiri_sparse_2022, jain_weighted_2020, chikhi_representation_2014}.\\
\textbf{The SSHash algorithm} (Sparse and Skew hashing of K-mers) introduced by \Citet{pibiri_sparse_2022} exploits these two properties by computing all minimizers appearing in a sequence and then use a minimal perfect hash function to store every super-k-mers offsets in an array \textit{Offsets} in the order induced by their minimizers hashes. In a second array \textit{Sizes} the $n$th element stores the starting index in \textit{Offsets} for the offsets belonging to the $n$th minimizer. The sparsity of minimizers mainly effects the data structures space consumption: By indexing super-k-mers, they can compress the data structures size by a factor of $(k-m+2)/2$ compared to indexing individual k-mers. The skewness of minimizers on the other hand is of importance to the locate time: Querying a k-mer $q$ requires computing its minimizer and then searching for it in every super-k-mer with the same minimizer. For a sufficiently large minimizer length, minimizers of high frequency become increasingly unlikely.\\
The both space consumption and locate time of SSHash are highly dependent of the minimizer length $m$: A smaller $m$ means a higher compression $(k-m+2)/2$. The lookup/locate time on the other hand is largest for small $m$ because those result in minimizers of higher \textit{frequencies}. This demonstrates the challenge that comes with choosing an optimal minimizer length.\\

\newcommand{\infrequent}{\textit{infrequent} }
\newcommand{\frequent}{\textit{frequent} }
\newcommand{\hashfunction}[1]{h_{#1}}
\textbf{The RSHash algorithm} (Rank-Select Hash) which was recently introduced by \citet{schulte-mattler_rshash_2026} is a minimizer based fast, compressed hash set for storing DNA k-mers that answers lookup and can be extended for supporting locate.\\
It follows a minimizer based approach like SSHash; however, instead of using minimal perfect hashing, RSHash relies on bitvectors supporting rank and select queries. Its main difference to previous methods is the introduction of a layered multi-minimizer scheme aiming to reduce the cost of high \textit{frequencies} whilst at also maintaining large compression.\\
The data structure works with one, two or three so called levels and a hash set. In order to build RSHash to answer lookup (or locate) queries one has to specify the number of levels as well as a pair of level parameters for each level $\level$: the minimizer length $\minimizerLength^\level$ and a threshold $\threshold^\level$.\\

\textbf{Building RSHash}: When using only one level, that level is build by first computing each k-mers minimizer of length $\minimizerLength^1$ and counting the number of times each minimizer occurs in the sequences. The minimizer computation uses a hash function $\hashfunction{1}$ to determine a total random order of the minimizers.\\
We call minimizers that appear at least once but no more than the threshold $\threshold^1$ to be \infrequent on the level.\\ 
A bitvector $R^1$ of lengths $4^{\minimizerLength^1}$ is used to mark the appearance of each distinct \infrequent minimizers. For this, the minimizers are interpreted as binary numbers of $2\minimizerLength^1$ bits and serve directly as indices to the vector. If for example the minimizers length were chosen to be $\minimizerLength^1=2$ and $AA$ was an \infrequent minimizer in the sequences, then $R^1[AA] = R^1[0]$ would be set to one.\\ 
The offsets each \infrequent minimizers occurrences in the sequences are stored in the offset vector $O^1$ in the order induced by $R^1$. Let $f$ for example be the frequency of the smallest \infrequent minimizer, that is the number of occurrences of the smallest minimizer for which the bit in $R^1$ is set. Then $O^1[0:n)$ contains the offsets of that minimizer whilst the next offsets $O^1[n:...)$ belong to the second smallest minimizer and so forth.\\
The boundaries of each minimizers offsets in $O^1$ are marked in the bitvector $S^1$ of length $|O^1|+1$. Following up on the example above, the first bit $S^1[0]$ would be set to one, indicating the start of the offsets belonging to the smallest minimizer, this would be followed by zeros until the $n$th bit $S^1[n] = 1$ indicating the start of the offsets belonging to the second smallest minimizer and so forth. The last bit in $S$ is also set to one marking the end of the largest minimizers offsets.\\

If a second level is specified, only those super-k-mers with \frequent minimizers on the first level are handled on the second level which is build almost identically to the first level, creating the vectors $R^2, O^2$ and $S^2$ but using the hash function $\hashfunction{2}$ for the computation of minimizers. The use of a different hash function for the computation leads to a skewed distribution of minimizers, in which most k-mers having \frequent minimizers on the first level, have \infrequent ones on the second. This holds even when the specified minimizer length $\minimizerLength^2$ and threshold $\threshold^2$ are equal to those on the first.\\
A third level would be build analogously, if specified.\\

Only the appearance of k-mers that have \frequent minimizers on all specified levels are then directly stored in a fast hash set sacrificing some space in return for a fast lookup which is viable as those k-mers are only a small fraction of the whole input sequence for a good choice of the level parameters.\\

\textbf{However, the choice of "good" parameters for a specific genome remains an open problem.} We already noted that using a small $m \ll k$ leads to a better compression but at the same time results in longer query times. This inverted effect of the minimizer length on the space consumption and query time is an inherent problem of minimizer based k-mer indexing methods.\\
Existing methods use a number of approaches to deal with this problem. For example, \citet{pibiri_sparse_2022} chooses a relatively large minimizer length, accepting a smaller compression in order to maintain reasonable locate times.\\

The leveled approach introduced in the RSHash algorithm on the other hand aims to achieve high compression by using relatively short minimizers on the lower levels whilst at the same time mitigating the burden of very \frequent minimizers by using thresholds to cut them off and guaranty fast queries on every level. The use of distinct hash function to compute minimizers on each level leads to different skewed distributions of the minimizers on each level. Because only a small percentage of the k-mers are frequent for each distribution, the number of k-mers reduces drastically with each level and only a small percentage of the sequences k-mers end up being stored in the spatially expensive hash set.\\

RSHash thus allow significantly faster queries opposed to SSHash whilst remaining similarly compact.
However, this is highly dependent of both the chosen minimizer lengths which influence the sparsity and skewness of minimizer distributions on each level and of the chosen thresholds, which affect both the lookup time on each level as well as the percentage of k-mers having to be stored in the hash set. For example, choosing very short minimizers and small thresholds on every level would result in a very large hash set as almost no minimizer would be \infrequent.\\

Choosing "good" level parameters to build the data structure for a given set of sequences currently requires manual testing and extensive understanding of the RSHash algorithm. 
We attempt to model the space consumption and lookup time of RSHash depending on the level parameters aiming to expand the understanding about how those values are influenced by the parameters and to work towards automating the choice of optimal level parameters for a given genome and optimization goal.

\chapter{Notation}\label{chap:notation}
%\newcommand{\superKmer}{sk}\todo{define super-k-mer}
\newcommand{\numberOfLevels}{{L}}
\newcommand{\setOfLevels}{\mathcal{L}}

\newcommand{\levelParameters}[1]{p^{#1}}
\newcommand{\varLevelParameters}[1]{\textbf{p}^{#1}}
\newcommand{\parameterAccess}[3]{\levelParameters{#1}[#2^{#3}]}

\newcommand{\frequencyOf}[1]{freq(#1)}
\newcommand{\minimizerOf}[2]{min^{\hashfunction{#2}}_{\minimizerLength^{#2}}(#1)}

\newcommand{\sequences}[1]{S_{\levelParameters{#1}}}
\newcommand{\substring}[3]{#1[#2:#3)}
A sequence is a string over the DNA alphabet $\Sigma = \{A,C,G,T\}$. Let $i,j \in \NN$ with $i < j$ be zero based indices of a string $S$ of length $|S| \geq j$. We use $\substring{S}{i}{j}$ to denote the substring from the $i$th to the $(j-1)$th character of $S$.
The k-mers of a string $S$ with $|S|\geq k$ are all its substrings $S[i:i+k) \in \Sigma^k$ of length $k$ for $i \in \NN \wedge i \leq |S|-k$.\\
Let $\numberOfLevels \in \NN^+$ be the number of levels specified. We refer to the set of levels $\setOfLevels := \{ \level \in \NN^+ \mid \level \leq \numberOfLevels \}$. We use $\possibleMinimizerLengths \subseteq [1,k) \cap \NN$ and $\possibleThresholds \subset \NN^+$ to denote the sets of minimizer lengths and thresholds considered. We denote the Cartesian $n$th power of a set $X$ with $X^n$ e.g. $X^3 = X \times X \times X$. A chosen minimizer length and threshold on a specified level $\level \in \setOfLevels$ are denoted $\minimizerLength^\level$ and $\threshold^\level$.\\
We use the tuple $\levelParameters{\level} := (\minimizerLength^1, ..., \minimizerLength^\level, \threshold^1, ..., \threshold^\level) \in \possibleMinimizerLengths^\level \times \possibleThresholds^\level$ to denote the level parameters up to level $\level$. For a specified tuple $\levelParameters{\level}$ of level parameters we use $\levelParameters{\level-1}$ to denote the level parameters $(\minimizerLength^1, ..., \minimizerLength^{\level-1}, \threshold^1, ..., \threshold^{\level-1})$ up to the previous level. Furthermore, we introduce the operator $\parameterAccess{\level}{\minimizerLength}{i} = \minimizerLength^i$ to denote retrieving the minimizer length specified for level $i$ from the level parameters and analogously $\parameterAccess{\level}{\threshold}{i}  = \threshold^i$.\\
We write $\sequences{\level-1} := \{S^\level_1,S^\level_2,...,S^\level_n\}$ to denote the set of sequences on level $\level$ in dependence of the level parameters $\levelParameters{\level-1}$ on the previous levels i.e. the k-mers which have \infrequent minimizers on all previous levels. The full set of input sequences when building RSHash do not depend on any parameters, are denoted with $\sequences{0}$ and are sometimes referred to as the input "text".\\
Let $\kmer \in \Sigma^{k}$ be a k-mer. The computation of its minimizer $\minimizerOf{\kmer}{\level}$ on level $l$ depends both on the specified minimizer length $\minimizerLength^\level$ as well as the hash function $\hashfunction{\level}$ on that level, which is not dependent of the level parameters but unique for every level. We denote the minimizers frequency with $\frequencyOf{\minimizerOf{\kmer}{\level}}$ which tells the combined number of the minimizer occurrences in a given set of input sequences.\\
We call the minimizer of $\kmer$ \infrequent on level $\level$ only if $1 \leq \frequencyOf{\minimizerOf{\kmer}{\level}} \leq \threshold^\level$ holds i.e. the minimizer occurs at least once but no more than the specified threshold on level $\level$. Minimizers with frequencies larger than the threshold are referred to as being \frequent.

\chapter{Sequence Statistics}\label{chap:sequenceStatistics}
As outlined in the \cref{chap:introduction}, the space consumption and query time of RSHash are highly dependent on the choice of level parameters as those influence the distributions of minimizers on all levels. We define properties of the minimizer distribution which will be needed for modeling the space consumption and lookup time in \cref{chap:minlp}. Specifically relevant for this is the behavior of minimizer distributions across several levels, which is unique to the approach proposed by \citet{schulte-mattler_rshash_2026} and therefore not considered in previous works.

\section{Minimizer Distribution} \label{sec:minimizerDistribution_severalLevels}
\newcommand{\distinctMinimizers}[1]{\mathcal{M}_{\levelParameters{#1-1}, \minimizerLength^{#1}}}
\newcommand{\distinctKmers}[1]{\mathcal{K}_{\levelParameters{#1-1}}}

\newcommand{\compactZV}[2]{#1[#2]}
\newcommand{\manualZV}[3]{#1[#2, #3]}
\newcommand{\ZVM}[1]{\manualZV{M}{\levelParameters{#1-1}}{\parameterAccess{#1}{\minimizerLength}{#1}}}
\newcommand{\ZVK}[1]{\manualZV{K}{\levelParameters{#1-1}}{\parameterAccess{#1}{\minimizerLength}{#1}}}

\newcommand{\probabilityDistribution}[1]{p_{#1}}

\newcommand{\CDF}[1]{F_{#1}}
\newcommand{\varCDF}[1]{\textbf{F}_{#1}}

Recall that we use $\levelParameters{\level} \in \possibleMinimizerLengths^\level \times \possibleThresholds^\level$ to describe the level parameters on all levels up to and including the level $\level$ and that $\sequences{0}$ is the full set of input sequences on the first level, which does not depend on any parameters. Also recall that the set $\sequences{\level-1}$ of input sequences on level $\level$ depends on the level parameters $\levelParameters{\level-1}$ on all previous level, as those affect, whether k-mers minimizer is \frequent or \infrequent.\\

Let $\distinctKmers{\level}$ be the set of distinct k-mers in a set of sequences $\sequences{\level-1}$. In other words, the set $\distinctKmers{\level}$ contains every distinct substring of length $k$ which appears in at least one of the sequences $S^\level_i \in \sequences{\level-1}$ which depend on the level parameters on previous levels:
\begin{equation}
	\distinctKmers{\level} := \{\kmer \in \Sigma^k \mid \exists S' \in \sequences{\level-1} ~ \exists (1 \leq i \leq |S'|-k) ~ S'[i:i+k) = \kmer \}
\end{equation}

Let $\distinctMinimizers{\level}$ be the set of distinct minimizers of lengths $\minimizerLength^\level$ on level $\level$ of the k-mers in $\distinctKmers{\level}$. The set of distinct minimizers is not only affected by the choice of level parameters on previous levels but also depends on the minimizer length on the current level:
\begin{equation}
	\distinctMinimizers{\level} := \{\minimizerOf{\kmer}{\level} \mid \forall \kmer \in \distinctKmers{\level}\}
\end{equation}
Let $\manualZV{M}{\levelParameters{\level-1}}{m}$ be a discrete random variable that describes the frequency of these minimizers:
\begin{align}
	\manualZV{M}{\levelParameters{\level-1}}{m} &: \distinctMinimizers{\level} \rightarrow \NN
	\\
	\manualZV{M}{\levelParameters{\level-1}}{m}(\minimizer) &= \frequencyOf{\minimizer}
\end{align}
We define its probability distribution $\probabilityDistribution{\manualZV{M}{\levelParameters{\level-1}}{\minimizerLength^\level}}$\footnote{This notation differs slightly from that used by \citet{schulte-mattler_rshash_2026} who do not introduce the probability distributions but instead use $p_M$ and $p_K$ to denote the cumulative distribution functions which we denote with $F_M$ and $F_K$}, which tells the probability of a minimizer having an exact frequency $f \in \NN$. The distribution of minimizers depends both on the level parameters $\levelParameters{\level-1}$ on previous levels, which affect the remaining k-mers on level $\l$ and the minimizer length $\minimizerLength^\level$ on the current level, which directly affects the computation of minimizers. We thus index the distributions with those parameters:
\begin{align}
	\probabilityDistribution{\manualZV{M}{\levelParameters{\level-1}}{\minimizerLength^\level}} &: \NN \rightarrow [0,1]
	\\
	\probabilityDistribution{\manualZV{M}{\levelParameters{\level-1}}{\minimizerLength^\level}}(f) &= \mathbb{P}(\manualZV{M}{\levelParameters{\level-1}}{\minimizerLength^\level} = f) = \frac{|\{\minimizer \in \distinctMinimizers{\level} \mid 
		\manualZV{M}{\levelParameters{\level-1}}{\minimizerLength^\level}(\minimizer)=f \}|}{|\distinctMinimizers{\level}|}
\end{align}
Using the random variables probability function we can derive its cumulative distribution function $\CDF{\manualZV{M}{\levelParameters{\level-1}}{m}}$, which returns the probability of a minimizer having a frequency that does not exceed an upper bound $f$: 
\begin{align}
	\CDF{\manualZV{M}{\levelParameters{\level-1}}{\minimizerLength^\level}} &: \NN \rightarrow [0,1]\\
	\CDF{\manualZV{M}{\levelParameters{\level-1}}{\minimizerLength^\level}}(f) &= \mathbb{P}(\manualZV{M}{\levelParameters{\level-1}}{\minimizerLength^\level} \leq f) = \sum_{i=1}^f \probabilityDistribution{\manualZV{M}{\levelParameters{\level-1}}{\minimizerLength^\level}}(i)
\end{align}

Whilst $\manualZV{M}{\levelParameters{\level-1}}{\minimizerLength^\level}$ and its distribution are useful for modeling values depending of a random (appearing) minimizer, RSHash is used to lookup k-mers, not minimizers. Therefore, we introduce the discrete random variable $\manualZV{K}{\levelParameters{\level-1}}{\minimizerLength^\level}$ which is defined for every distinct (appearing) k-mer on the level and describes the frequency each k-mers minimizer dependent of the level parameters $\levelParameters{\level-1}$ on all previous levels and the minimizer length $\minimizerLength^\level$ on the current one:
\begin{align}
	\manualZV{K}{\levelParameters{\level-1}}{\minimizerLength^\level} &: \distinctKmers{\level} \rightarrow \NN
	\\
	\manualZV{K}{\levelParameters{\level-1}}{\minimizerLength^\level}(\kmer) &= \frequencyOf{\minimizerOf{\kmer}{\level}}
\end{align}
Its probability distribution $\probabilityDistribution{\manualZV{K}{\levelParameters{\level-1}}{\minimizerLength^\level}}$ and its cumulative distribution function $\CDF{\manualZV{K}{\levelParameters{\level-1}}{\minimizerLength^\level}}$ are: 
\begin{align}
	\probabilityDistribution{\manualZV{K}{\levelParameters{\level-1}}{\minimizerLength^\level}} &: \NN \rightarrow [0,1]
	\\
	\probabilityDistribution{\manualZV{K}{\levelParameters{\level-1}}{\minimizerLength^\level}}(f) &= \mathbb{P}(\manualZV{K}{\levelParameters{\level-1}}{\minimizerLength^\level} = f) = \frac{|\{\kmer \in \distinctKmers{\level} \mid \manualZV{K}{\levelParameters{\level-1}}{\minimizerLength^\level}(\kmer)=f \}|}{|\distinctKmers{\level}|}
	\\
	\CDF{\manualZV{K}{\levelParameters{\level-1}}{\minimizerLength^\level}} &: \NN \rightarrow [0,1]\\
	\CDF{\manualZV{K}{\levelParameters{\level-1}}{\minimizerLength^\level}}(f) &= \mathbb{P}(\manualZV{K}{\levelParameters{\level-1}}{\minimizerLength^\level} \leq f) = \sum_{i=1}^f \probabilityDistribution{\manualZV{K}{\levelParameters{\level-1}}{\minimizerLength^\level}}(i)
\end{align}

Let $\levelParameters{\level}$ be a combination of level parameters on a level $\level$. Recall that we use $\parameterAccess{\level}{\minimizerLength}{\level}$ to denote retrieving the minimizer length specified for level $\level$ from the level parameters and analogously for $\parameterAccess{\level}{\threshold}{\level}$. For better readability we introduce a more compact notation to denote the probability $\CDF{\compactZV{M}{\levelParameters{\level}}}$ of a minimizer being \infrequent on level $\level$ and $\CDF{\compactZV{K}{\levelParameters{\level}}}$ of a k-mer having an \infrequent minimizer on that level, both in dependence of all level parameters:
\begin{align}
	\CDF{\compactZV{M}{\levelParameters{\level}}} &= \CDF{\manualZV{M}{\levelParameters{\level-1}}{\parameterAccess{\level}{\minimizerLength}{\level}}}(\parameterAccess{\level}{\threshold}{\level})\\
	\CDF{\compactZV{K}{\levelParameters{\level}}} &= \CDF{\manualZV{K}{\levelParameters{\level-1}}{\parameterAccess{\level}{\minimizerLength}{\level}}}(\parameterAccess{\level}{\threshold}{\level})
\end{align}

\section{Average Frequencies} \label{sec:averagefrequencies}
Let $\Omega_{\manualZV{M}{\levelParameters{\level-1}}{\minimizerLength^\level}}$ be the sample space of the random variable $\manualZV{M}{\levelParameters{\level-1}}{\minimizerLength^\level}$ i.e. the set which contains all frequencies $f$ for which there is at least one minimizer $\minimizer$ on level $\level$ which occurs exactly $f$ times:
\begin{equation}
	\Omega_{\manualZV{M}{\levelParameters{\level-1}}{\minimizerLength^\level}} = \{f \mid \exists \minimizer \in \distinctMinimizers{\level} \frequencyOf{\minimizer} = f \}
\end{equation}
The average frequency of minimizers on a level $\level \in \setOfLevels$ is given by the expected value of the probability distribution of the random variable $\manualZV{M}{\levelParameters{\level-1}}{m}$:
\begin{equation}
	\mathbb{E}[\manualZV{M}{\levelParameters{\level-1}}{\minimizerLength^\level}] = \sum_{f \in \Omega_{\manualZV{M}{\levelParameters{\level-1}}{\minimizerLength^\level}}} \probabilityDistribution{\manualZV{M}{\levelParameters{\level-1}}{\minimizerLength^\level}}(f) * f
\end{equation}
For a given set of finite sequences the set of distinct minimizers on each level and for each combination of level parameters is finite by definition. It follows that for a given sequence and set of level parameters, the average frequency $\mathbb{E}[\manualZV{M}{\levelParameters{\level-1}}{\minimizerLength^\level}]$ can always be computed, because the sample space $\Omega_{\manualZV{M}{\levelParameters{\level-1}}{\minimizerLength^\level}}$ is then finite by definition.\\

More relevant to the optimization of level parameters is the average frequency of \infrequent minimizers i.e. the average frequency of minimizers whose frequency do not exceed the threshold $\threshold^\level$. We denote this value with $\averageFrequency(\levelParameters{\level})$ and can derive it from $\probabilityDistribution{\manualZV{M}{\levelParameters{\level-1}}{m}}$ by computing a weighted sum over all frequencies up until the threshold $\threshold^\level$ and then norming the result by dividing through the percentage of \infrequent minimizers:
\begin{equation} \label{equation:averageFrequency}
	\averageFrequency(\levelParameters{\level}) = \frac{1}{\CDF{\compactZV{M}{\levelParameters{\level}}}} \cdot
	\sum\limits_{f=1}^{\threshold^\level} \probabilityDistribution{\manualZV{M}{\levelParameters{\level-1}}{\parameterAccess{\level}{\minimizerLength}{\level}}}(f) \cdot f 
\end{equation}

\section{Numbers of Distinct \textit{Infrequent} Minimizers}
\newcommand{\numberOfDistinctMinimizers}[1]{n_{DIM}(#1)}
\newcommand{\varNumberOfDistinctMinimizers}{\textbf{n}_{\textbf{DIM}}}
We define the number of distinct \infrequent minimizers $\numberOfDistinctMinimizers{\levelParameters{\level}}$ which also depend on all level parameters $\levelParameters{\level}$ of a level $\level$.  Their value can be computed py multiplying the total number of distinct minimizers on a level with the percentage of \infrequent ones:
\begin{equation}
	\numberOfDistinctMinimizers{\levelParameters{\level}} = |\distinctMinimizers{\level}| \cdot \CDF{\compactZV{M}{\levelParameters{\level}}}
\end{equation}


\chapter{Mixed Integer Non Linear Program}\label{chap:minlp}
We model the optimization problem as a Mixed Integer Non Linear Programming (MINLP). In \cref{sec:minlp_oneLevel} we model the lookup time and space needed by RSHash on one level in dependence of the level parameters and minimizers distribution. We then extend this to the case of several levels in \cref{sec:minlp_severalLevels}.\\
In \cref{chap:experiments} we describe how to find the optimal level parameters for a given sequence using the MINLP.

\section{Model Input}\label{sec:minlp_modelInput}
\newcommand{\nkmers}{n_{kmers}}
The set of levels $\setOfLevels$, the sets of minimizer lengths $\possibleMinimizerLengths \subseteq [1,k) \cap \NN$ and thresholds $\possibleThresholds \subset \NN$ to consider have to be specified explicitly.
\\
Furthermore, the number $\nkmers$ of k-mers in the sequences the values $\CDF{\compactZV{M}{\levelParameters{\level}}}$, $\CDF{\compactZV{K}{\levelParameters{\level}}}$,  $\averageFrequency(\levelParameters{\level})$ and $\numberOfDistinctMinimizers{\levelParameters{\level}}$ for all possible combinations $\levelParameters{\level}$ of the level parameters on each level $\level \in \setOfLevels$ as defined in \cref{chap:sequenceStatistics} have to be passed as input to the MINLP. Their computation is further explained in \cref{chap:experiments}. 
\\
Lastly, a measure of the times 
$\timeForOperationOnStructure{\contains}{R}$,
$\timeForOperationOnStructure{\select}{S}$,
$\timeForOperationOnStructure{\access}{O}$,
$\timeForOperationOnStructure{\compare}{T}$,
$\timeForOperationOnStructure{check}{H}$
needed for a contains rank query on bitvector $R$, a select queries on bitvector $S$, accessing an offset in $O$, a comparison with the text $T$ at a random offset and checking whether a key exists in the hash set $H$. The computation of the minimizer distribution as well as the times needed to perform these actions on the specified data structures is outlined in \cref{chap:experiments}.

\section{Model for One Level}\label{sec:minlp_oneLevel}
Aiming to improve readability, we omit the superscript index indicating an instances belonging to a specific level as only one level is considered in this section e.g. the minimizer length is simply referred to as $\minimizerLength$ instead of $\minimizerLength^1$.
\\
In order to determine $\varMinimizerLength$ and $\varThreshold$ we use binary indicator variables $\indicatorm_\minimizerLength$ and $\indicatort_\threshold$ for each possible $\minimizerLength \in \possibleMinimizerLengths$ and $\threshold \in \possibleThresholds$:
\begin{align}
	\indicatorm_\minimizerLength &\in \{0,1\} \mid \forall \minimizerLength \in \possibleMinimizerLengths \label{variable:indicatorsM}\\
	\indicatort_\threshold &\in \{0,1\} \mid \forall \threshold \in \possibleThresholds \label{variable:indicatorsT}
\end{align} 

The indicator variables for each minimizer length should be mutually exclusive as $\varMinimizerLength$ can only take one distinct value. The same holds for the indicator variables for each threshold:
\newcommand{\mututallyExclusiveM}{\sum\limits_{\minimizerLength \in \possibleMinimizerLengths} \indicatorm_\minimizerLength &= 1}
\newcommand{\mututallyExclusiveT}{\sum\limits_{\threshold \in \possibleThresholds} \indicatort_\threshold &= 1}
\begin{align}
	\mututallyExclusiveM \label{constraint:mututallyExclusiveM}\\
	\mututallyExclusiveT \label{constraint:mututallyExclusiveT}
\end{align}

The variables $\varMinimizerLength \in \possibleMinimizerLengths$ and $\varThreshold \in \possibleThresholds$ are introduced to describe the optimal minimizer length and threshold.\\
There values can be modeled through a sum over the indicator variables multiplied with there corresponding values:
\newcommand{\optimalM}{\varMinimizerLength &= \sum\limits_{\minimizerLength \in \possibleMinimizerLengths} \indicatorm_\minimizerLength \cdot \minimizerLength}
\newcommand{\optimalT}{\varThreshold &= \sum\limits_{\threshold \in \possibleThresholds} \indicatort_\threshold \cdot \threshold}
\begin{align}
	\varMinimizerLength &\in \possibleMinimizerLengths \label{variable:optimalM}\\
	\varThreshold &\in \possibleThresholds \label{variable:optimalT}\\
	\optimalM \label{constraint:optimalM}\\
	\optimalT \label{constraint:optimalT}
\end{align}

We introduce the variable $\varCDF{M} \in [0,1]$ to take the value of $\CDF{\compactZV{M}{(\varMinimizerLength, \varThreshold)}}$ i.e. the percentage of minimizers of length $\varMinimizerLength$ in the input sequence, whose frequencies do not exceed $\varThreshold$ (\cref{sec:minimizerDistribution_severalLevels}).\\
It is not possible to directly evaluate this value as both $\varMinimizerLength$ and $\varThreshold$ are variables of the model. Instead we use the indicator variables to model a constraint forcing $\varCDF{M}$ to take the value of $\CDF{\compactZV{M}{(\minimizerLength, \threshold)}}$ for which $\indicatorm_\minimizerLength = 1$ and $\indicatort_\threshold = 1$:
\newcommand{\cdfMExpression}{
	\varCDF{M} &= \sum\limits_{(\minimizerLength, \threshold) \in \possibleMinimizerLengths \times \possibleThresholds} \CDF{\compactZV{M}{(\minimizerLength, \threshold)}} \cdot
	\indicatorm_\minimizerLength \indicatort_\threshold
}
\begin{align}
	\varCDF{M} &\in [0,1] \label{variable:percMinimizers}\\
	\cdfMExpression \label{constraint:percMinimizers}
\end{align}

Analogously we introduce $\varCDF{K} \in [0,1]$ to represent $\CDF{\compactZV{K}{(\varMinimizerLength, \varThreshold)}}$ and $\varAverageFrequency \in \RR^+_0$ to represent $\averageFrequency((\varMinimizerLength, \varThreshold))$ (\cref{sec:averagefrequencies}):
\newcommand{\cdfKExpression}{
	\varCDF{K} &= \sum\limits_{(\minimizerLength, \threshold) \in \possibleMinimizerLengths \times \possibleThresholds} \CDF{\compactZV{K}{(\minimizerLength, \threshold)}} \cdot
	\indicatorm_\minimizerLength \indicatort_\threshold
}
\newcommand{\varAverageFrequencyExpression}{
	\varAverageFrequency &= \sum\limits_{(\minimizerLength, \threshold) \in \possibleMinimizerLengths \times \possibleThresholds} \averageFrequency((\minimizerLength,\threshold)) \cdot
	\indicatorm_\minimizerLength \indicatort_\threshold
}
\begin{align}
	\varCDF{K} &\in [0,1] \label{variable:percKmers}\\
	\varAverageFrequency &\in \RR^+_0 \label{variable:averageFrequency}\\
	\cdfKExpression \label{constraint:percKmers}\\
	\varAverageFrequencyExpression \label{constraint:averageFrequency}
\end{align}

\subsection{Lookup Time}
In order to model the time needed for a lookup using RSHash, we closely follow the lookup procedure as described by \citet{schulte-mattler_rshash_2026}:\\
A k-mer $\kmer$ is looked up by first computing its minimizer $\minimizerOf{\kmer}{}$ and checking whether the bit $R[\minimizerOf{\kmer}{}]$ is set whilst simultaneously computing its rank $r = \rank_R(\minimizerOf{\kmer}{})$, taking $\timeForOperationOnStructure{\contains}{R}$ time (part of the model input \cref{sec:minlp_modelInput}).\\
For the percentage $\varCDF{K}$ (\cref{constraint:percKmers}) of the k-mers this bit is set and the lookup is completed by searching for $\kmer$ on the level. We introduce the variable $\timeForOperationOnStructure{level}{} \in \RR^+_0$ to describe the time needed for this.\\
If the bit $R[\minimizerOf{\kmer}{}]$ is not set, the lookup is completed through a lookup in the hash set $H$ taking $\timeForOperationOnStructure{check}{H}$ time (part of the model input \cref{sec:minlp_modelInput}). The variable $\timeForOperationOnStructure{lookup}{} \in \RR^+_0$ describing the lookup time can thus be modeled as:
\newcommand{\lookupTimeExpression}{
	\nonumber \timeForOperationOnStructure{lookup}{} =&~ \timeForOperationOnStructure{\contains}{R} + \varCDF{K} \cdot \timeForOperationOnStructure{level}{}\\
	&~\phantom{\timeForOperationOnStructure{\contains}{R}}+ (1-\varCDF{K}) \cdot \timeForOperationOnStructure{check}{H}}
\begin{align}
	\timeForOperationOnStructure{level}{} \in&~ \RR^+_0 \label{variable:levelTime}\\
	\timeForOperationOnStructure{lookup}{} \in&~ \RR^+_0 \label{variable:lookupTime}\\
	\lookupTimeExpression \label{constraint:lookupTime}
\end{align}

The search for those k-mers $\kmer$ for which the bit $R[\minimizerOf{\kmer}{}]$ is set continues on the level by using the minimizers rank  $r = \rank_R(\minimizerOf{\kmer}{})$ in the select\todo{explain select and rank} queries $first = \select_S(r)$ and $last = \select_S(r+1)-1$ which tell the indices in Offsets $O$ where the minimizer’s first and last offset in the text $T$ are stored. These two queries take $2\timeForOperationOnStructure{\select}{S}$ time.\\
Finally the presence of $\kmer$ in $T$ is checked at every minimizer offset in the interval $O[first:last]$ until $\kmer$ is either found or the end of the interval is reached. The average frequency $\varAverageFrequency$ of \infrequent minimizers is exactly the average length of this interval. In the lookup case, the comparison is terminated as soon as $\kmer$ is found for the first time, which is on average the case after $(\varAverageFrequency+1)/2$ comparisons with each comparison taking $\timeForOperationOnStructure{\compare}{T}$ time.\\
The following constraint models the complete time $\timeForOperationOnStructure{level}{}$ needed to search for a k-mers presence on the level:
\newcommand{\levelTime}{\timeForOperationOnStructure{level}{} &= 2\timeForOperationOnStructure{\select}{S} + (\varAverageFrequency+1)/2 \cdot (\timeForOperationOnStructure{\access}{O} + \timeForOperationOnStructure{\compare}{T})}
\begin{align}
	\levelTime \label{constraint:levelTime}
\end{align}
Note that the value $\timeForOperationOnStructure{lookup}{}$ is the average time for a positive RSHash lookup, that is the time to lookup a k-mer which was randomly sampled from the sequence. In practice, queried k-mers may or may not appear in the sequence - hence the need for a lookup. Looking up a k-mer which does not appear in the sequence (negative lookup) would always take $\timeForOperationOnStructure{\contains}{R} + \timeForOperationOnStructure{check}{H}$; however, we do not account for this as the ratio of positive and negative queries is impossible predict.\\

Without further constraints, solving the described MINLP for minimal $\timeForOperationOnStructure{lookup}{}$ results in a solution where the $\indicatorm_\minimizerLength$, $\indicatort_\threshold$ corresponding to the smallest $m \in \possibleMinimizerLengths$, $t \in \possibleThresholds$ will be set to one.\\
This is expected, as this leads to the smallest possible $\varCDF{K}$ which minimizes the k-mers having to be looked up in the text and maximizes the lookups in the (faster) hash set.\\
However, this results in maximal space consumption, as each entry in the hash set is significantly larger than an entry in $O$ as described in \cref{sec:problemDefinition}.\\
In order to account for that, we have to add constraints modeling the space consumption.

\subsection{Space Consumption}
Instead of modeling the total space used by RSHash we will focus on the bits per k-mer (bpk) used by the data structure as this metric allows to compare compression independent of the sequence length.\\
The space consumption of RSHash can be decomposed into the space needed for the data structures components: the text $T$\footnote{Recall that we use the 'text' $T$ to refer to the set of input sequences $\sequences{0}$ in order to avoid confusion with the bit vector $S$.}, the vectors $R$, $S$ and $O$ as well as the hash set $H$. We use one variable to describe the bits per k-mer needed by each of the components:
\begin{align}
	bpk_T &\in \RR_0^+ \label{variable:bpkT}\\
	bpk_R &\in \RR_0^+ \label{variable:bpkR}\\
	bpk_S &\in \RR_0^+ \label{variable:bpkS}\\
	bpk_O &\in \RR_0^+ \label{variable:bpkO}\\
	bpk_H &\in \RR_0^+ \label{variable:bpkH}
\end{align}

\textbf{The text $T$} is much longer than k. Because of this the bits per k-mer for storing the text $T$ can be set to a constant of $2$ bits needed to store one base:
\newcommand{\bpkT}{bpk_T &= 2}
\begin{align}
	\bpkT \label{constraint:bpkT}
\end{align} 
 
\textbf{Bitvector $R$} is a bitvector of length $4^\varMinimizerLength$ in which each bit set to $1$ marks the occurrence of one distinct \infrequent minimizer in the sequences. The amount of ones in $R$ can thus be computed by selecting the the With $\numberOfDistinctMinimizers{\levelParameters{(\minimizerLength, \threshold)}}$ being the number of distinct \infrequent minimizers in dependence of the specified level parameters, we can derive the value $\varNumberOfDistinctMinimizers$ for the chosen minimizer length $\varMinimizerLength$ and threshold $\varThreshold$:
\newcommand{\nDIMExpression}{
	\varNumberOfDistinctMinimizers &= \sum\limits_{(\minimizerLength, \threshold) \in \possibleMinimizerLengths \times \possibleThresholds}
	\numberOfDistinctMinimizers{(\minimizerLength, \threshold)} \cdot
	\indicatorm_\minimizerLength \indicatort_\threshold
}
\begin{align}
	\varNumberOfDistinctMinimizers &\in \RR^+_0 \label{variable:nDIM}\\
	\nDIMExpression \label{constraint:nDIM}
\end{align}
We note that for continuous sequences one could also model $\varNumberOfDistinctMinimizers$ by computing the number $\varCDF{K} \cdot \nkmers$ of k-mers with \infrequent minimizers and dividing it by the average frequency $\varAverageFrequency$ of minimizers multiplied with the ratio $(k-\varMinimizerLength+2)/2$ of k-mers covered per minimizer. This should decrease the models complexity as it would replace the non-linear constraint $\cref{constraint:nDIM}$ with a linear one. However, the ratio $(k-\varMinimizerLength+2)/2$ is only valid for continuous sequences $S$ of lengths $|S| \gg k > \varMinimizerLength$, rendering it useless both for sequences which have been preprocessed to remove duplicate k-mers and for \cref{sec:minlp_severalLevels} where the input on higher levels which consists of many relatively short super-k-mers.\\

Using Elias-Fano-Encoding, a bitvector of size $U$ with $n$ bits set to one can be compressed to
$n (\lceil log_2(U/n) \rceil + 3)$ \cite{okanohara2007practical}. Therefore, we can model the bits per k-mer needed to store $R$ by using this compression and dividing this by the total number $\nkmers$ of k-mers yields the constraint describing the bits per k-mer of $R$:
\newcommand{\bpkR}{bpk_{R} &= \frac{\varNumberOfDistinctMinimizers}{\nkmers} \cdot \left(\lceil 2\varMinimizerLength - log_2\varNumberOfDistinctMinimizers \rceil + 3 \right)}
\begin{align}
	\bpkR \label{constraint:bpkR}
\end{align}
%\todo[inline]{ start experimental}
%\textbf{Bit Vector Rank Query}
%A bitvector of length $U$ with $n$ $1s$ can be encoded with $n\lceil log_2 U/n \rceil + 3n$ many bits supporting rank in $O(log_2(U/n)$ time \cite{okanohara2007practical}.
%\begin{equation*}
%	\timeForOperationOnStructure{\contains}{R} = log_2 \frac{4^\varMinimizerLength}{\varNumberOfDistinctMinimizers} =  2\varMinimizerLength - log_2\varNumberOfDistinctMinimizers
%\end{equation*}
%\todo[inline]{end experimental}

\textbf{Bitvector $S$} has length $|O|+1 \approx |O|$ where every bit corresponds to the offset of one minimizer encountered in the sequences and one extra bit marks the end of all offsets - this last bit can be disregarded when modeling $bpk_S$ as its effect is insignificantly small.\\
Multiplying the number $\varNumberOfDistinctMinimizers$ of distinct \infrequent minimizers with their average frequency $\varAverageFrequency$ gives the length of the offset vector $O$. To model the bits per k-mer in $S$ we divide this by the number of k-mers in the sequence:
\newcommand{\bpkS}{bpk_{S}	&= \frac{\varNumberOfDistinctMinimizers \cdot \varAverageFrequency}{\nkmers}}
\begin{align}
	\bpkS \label{constraint:bpkS} 
\end{align}
Note that if one figured out, how to model the ratio between minimizers and k-mers on non-continuous sequences, $bpk_S$ could be modeled without the use of $\varNumberOfDistinctMinimizers$ by taking the percentage $\varCDF{K}$ of k-mers with infrequent minimizers and dividing it by that ratio.
%\todo[inline]{ start experimental}
%\textbf{Bit Vector Rank Query}
%A bitvector of length $U$ with $n$ $1s$ can be encoded with $n\lceil log_2 U/n \rceil + 3n$ many bits supporting rank in $O(log_2(U/n)$ time \cite{okanohara2007practical}.
%\begin{equation*}
%	\timeForOperationOnStructure{\select}{S} = log_2 \frac{\varCDF{M} \cdot n_{minimizers}}{\varNumberOfDistinctMinimizers} = \varAverageFrequency
%\end{equation*}
%\todo[inline]{end experimental}

\textbf{The offset vector $O$} contains one entry of size $log_2 |T|$ for the offset of each \infrequent minimizer.\\
Since $bpk_S$ already describes the number of offsets per k-mer, we can model the bits per k-mer in $O$ with:
\newcommand{\bpkO}{bpk_O &= bpk_S \cdot log_2 |T|}
\begin{align}
	\bpkO \label{constraint:bpkO}
\end{align}

\textbf{The hash set $H$} contains an entry $(key:\kmer)$ for each distinct k-mer $\kmer$ with an \infrequent minimizer, which are $1-\varCDF{K}$ of all k-mers in $T$.
Each entry has size $size(key) = 2k$ bit. We introduce the load factor\footnote{In practice, the load factor of a hash set is usually not known before building the set but rather a value used for re scaling. Therefore $\alpha=0.75$ is rather an estimate than the actual load factor of a possible hash set. We further asses the consequences of this in \cref{chap:experiments}.} $\alpha=0.75$.
The bits per k-mer required for the hash set can thus be modeled by:
\newcommand{\bpkH}{bpk_H &= (1-\varCDF{K}) \cdot \frac{2k}{\alpha}}
\begin{align}
	\bpkH \label{constraint:bpkH}
\end{align}

\textbf{The total bits per k-mer} $bpk_{total} \in \RR^+_0$ can be modeled by simply adding up the space requirements of each component:
\newcommand{\bpkTotal}{bpk_{total} &= bpk_T + bpk_R + bpk_S + bpk_O + bpk_H}
\begin{align}
	bpk_{total} &\in \RR^+_0 \label{variable:bpkTotal}\\
	\bpkTotal \label{constraint:bpkTotal}
\end{align}

\subsection{Optimization Goal}
Having defined all constraints regarding the lookup time and space consumption of RSHash, we can now construct a MINLP in order to find optimal level parameters.\\
We introduce three more external parameters to the model: The optimization goal $OptGoal$ which specifies whether to minimize $\timeForOperationOnStructure{lookup}{}$ or $bpk_{total}$ and upper bounds $Bound_{time}, Bound_{bpk}$ for the lookup time and space consumption. 
\begin{align*}
	OptGoal			&\in \{\timeForOperationOnStructure{lookup}{}, bpk_{total}\} \\
	Bound_{time}	&\in \RR_0^+\\
	Bound_{bpk}		&\in \RR_0^+
\end{align*}
When using the MINLP, the upper bound belonging to the optimization goal should be set to $\infty$. In addition to the already defined ones, the following two constraints are added to the MINLP to enforce these bounds:
\begin{align}
	\timeForOperationOnStructure{lookup}{} &\leq Bound_{time} \label{constraint:boundTime}\\
	bpk_{total} &\leq Bound_{bpk} \label{constraint:boundBpk}
\end{align}
One could for example search for level parameters minimizing the lookup time whilst not exceeding eight bits per k-mer by setting $OptGoal \leftarrow \timeForOperationOnStructure{lookup}{}, Bound_{time} \leftarrow \infty, Bound_{bpk} \leftarrow 8$. The constructed MINLP for this example is:
\begin{align*}
	INPUT: ~
	&OptGoal = \timeForOperationOnStructure{lookup}{},
	Bound_{time} = \infty,
	Bound_{bpk} = 8,
	\\
	&\setOfLevels = \{1\}, \possibleMinimizerLengths, \possibleThresholds, \nkmers
	\\
	&\CDF{\compactZV{M}{(\minimizerLength, \threshold)}}, \CDF{\compactZV{K}{(\minimizerLength, \threshold)}},
	\averageFrequency((\minimizerLength, \threshold)),
	\numberOfDistinctMinimizers{(\minimizerLength, \threshold)} ~ \forall (\minimizerLength, \threshold) \in \possibleMinimizerLengths \times \possibleThresholds,
	\\
	&\timeForOperationOnStructure{\contains}{R},
	\timeForOperationOnStructure{\select}{S},
	\timeForOperationOnStructure{\access}{O},
	\timeForOperationOnStructure{\compare}{T},
	\timeForOperationOnStructure{check}{H}
\end{align*}

\begin{align*}
	\text{MINIMIZE}\\
		&OptGoal = \timeForOperationOnStructure{lookup}{}\\
	\text{SUBJECT TO}\\
	\mututallyExclusiveM &\cref{constraint:mututallyExclusiveM} \\
	\mututallyExclusiveT &\cref{constraint:mututallyExclusiveT} \\
	\optimalM &\cref{constraint:optimalM}\\
	\optimalT &\cref{constraint:optimalT}\\
	\cdfMExpression &\cref{constraint:percMinimizers}\\
	\cdfKExpression &\cref{constraint:percKmers}\\
	\varAverageFrequencyExpression &\cref{constraint:averageFrequency}\\
	\levelTime &\cref{constraint:levelTime}\\
	\lookupTimeExpression &\cref{constraint:lookupTime}\\
	\timeForOperationOnStructure{lookup}{} &\leq Bound_{time} = \infty &\cref{constraint:boundTime}\\
	\bpkT &\cref{constraint:bpkT}\\
	\nDIMExpression &\cref{constraint:nDIM}\\
	\bpkR &\cref{constraint:bpkR}\\
	\bpkS &\cref{constraint:bpkS}\\
	\bpkO &\cref{constraint:bpkO}\\
	\bpkH &\cref{constraint:bpkH}\\
	\bpkTotal &\cref{constraint:bpkTotal}\\
	bpk_{total} &\leq Bound_{bpk} = 8 &\cref{constraint:boundBpk}\\
	%		\varMinimizerLength &= \sum\limits_{m \in \possibleMinimizerLengths} \indicatorm_\minimizerLength \cdot m \\
	%		\varThreshold &= \sum\limits_{t \in \possibleThresholds} \indicatort_\threshold \cdot t
\end{align*}

\begin{align*}
	\text{VARIABLES}\\
	\indicatorm_\minimizerLength &\in \{0,1\} \mid \forall m \in M &\cref{variable:indicatorsM}\\
	\indicatort_\threshold &\in \{0,1\} \mid \forall t \in T &\cref{variable:indicatorsT}\\
	\varMinimizerLength &\in \possibleMinimizerLengths &\cref{variable:optimalM}\\
	\varThreshold &\in \possibleThresholds &\cref{variable:optimalT}\\
	\varCDF{M} &\in [0,1] &\cref{variable:percMinimizers}\\
	\varCDF{K} &\in [0,1] &\cref{variable:percKmers}\\
	\varAverageFrequency &\in \RR_0^+ &\cref{variable:averageFrequency}\\
	\timeForOperationOnStructure{level}{} &\in \RR^+_0 &\cref{variable:levelTime}\\
	\timeForOperationOnStructure{lookup}{} &\in \RR_0^+ &\cref{variable:lookupTime}\\
	bpk_T &\in \RR_0^+ &\cref{variable:bpkT}\\
	bpk_R &\in \RR_0^+ &\cref{variable:bpkR}\\
	bpk_S &\in \RR_0^+ &\cref{variable:bpkS}\\
	bpk_O &\in \RR_0^+ &\cref{variable:bpkO}\\
	bpk_H &\in \RR_0^+ &\cref{variable:bpkH}\\
	\varNumberOfDistinctMinimizers &\in \RR^+_0 &\cref{variable:nDIM}\\
	bpk_{total} &\in \RR_0^+ &\cref{variable:bpkTotal}
\end{align*} 

\section{Model for Several Levels}\label{sec:minlp_severalLevels}
In the following, we adapt the constraints defined for using only one level in \cref{sec:minlp_oneLevel} to model the use of $\numberOfLevels \in \NN^+$ levels with $\setOfLevels = \{ \level \in \NN^+ \mid \level \leq \numberOfLevels \}$ being the set of levels used\footnote{We model the space consumption and lookup time for any positive natural number of levels even though RSHash currently only supports the use of one, two or three levels \cite{schulte-mattler_rshash_2026}.}.\\
When modeling several levels, we add a superscript $\level$ to variables in order to indicate there belonging to a level $\level \in \setOfLevels$ e.g. $\varMinimizerLength^1, \varMinimizerLength^2$ and $\varMinimizerLength^3$ denote the optimal minimizer lengths on level $1,2$ and $3$.\\

Unique sets of binary indicator variables $\indicatorm^\level_\minimizerLength$ and $\indicatort^\level$ are needed for the level parameters $\varMinimizerLength^\level$ and $\varThreshold^\level$ of each level $\level$.
The constraints regarding these variables are analog to those when using only one level:
\[ \begin{array}{rll}
	\indicatorm_\minimizerLength^\level &\in \{0,1\} & \mid \forall (\level,\minimizerLength) \in (\setOfLevels \times \possibleMinimizerLengths) \\
	\indicatort_\threshold^\level &\in \{0,1\} & \mid \forall (\level,\threshold) \in (\setOfLevels \times \possibleThresholds) \\
	\sum\limits_{m \in \possibleMinimizerLengths} \indicatorm_\minimizerLength^\level &= 1 & \mid \forall l \in \setOfLevels\\
	\sum\limits_{t \in \possibleThresholds} \indicatort_\threshold^\level &= 1 & \mid \forall l \in \setOfLevels\\
	\varMinimizerLength^\level &\in \RR^+_0 & \mid \forall l \in \setOfLevels\\
	\varThreshold^\level &\in \RR^+_0 & \mid \forall l \in \setOfLevels\\
	\varMinimizerLength^\level &= \sum\limits_{m \in \possibleMinimizerLengths} \indicatorm_\minimizerLength^\level \cdot \minimizerLength^\level & \mid \forall l \in \setOfLevels \\
	\varThreshold^\level &= \sum\limits_{t \in \possibleThresholds} \indicatort_\threshold^\level \cdot \threshold^\level & \mid \forall l \in \setOfLevels
\end{array} \]
For each level $l$ we model $\varCDF{M}^\level$, $\varCDF{K}^\level$ and $\varAverageFrequency^\level$ analogously to using a single level; however, it is necessary to sum over every possible combination of level parameters $\levelParameters{\level}$. The number of additions thus grows exponentially for an increasing number of levels.
\begin{align*}
	\varCDF{M}^\level &\in [0,1] &\mid \forall l \in \setOfLevels
	\\
	\varCDF{K}^\level &\in [0,1] &\mid \forall l \in \setOfLevels
	\\
	\varAverageFrequency^\level &\in \RR^+_0 &\mid \forall l \in \setOfLevels
	\\
	\varCDF{M}^\level &= \sum\limits_{\levelParameters{\level} \in \possibleMinimizerLengths^\level \times \possibleThresholds^\level}
	\CDF{\compactZV{M}{\levelParameters{\level}}} 
	\cdot
	\prod_{1\leq i \leq \level} \indicatorm^i_{\parameterAccess{\level}{\minimizerLength}{i}} \indicatort^i_{\parameterAccess{\level}{\threshold}{i}}
	& \mid \forall l \in \setOfLevels
	\\
	\varCDF{K}^\level &= \sum\limits_{\levelParameters{\level} \in \possibleMinimizerLengths^\level \times \possibleThresholds^\level}
	\CDF{\compactZV{K}{\levelParameters{\level}}}
	\cdot
	\prod_{1\leq i \leq \level} \indicatorm^i_{\parameterAccess{\level}{\minimizerLength}{i}} \indicatort^i_{\parameterAccess{\level}{\threshold}{i}}
	& \mid \forall l \in \setOfLevels
	\\
	\varAverageFrequency^\level &= \sum\limits_{\levelParameters{\level} \in \possibleMinimizerLengths^\level \times \possibleThresholds^\level}
	\averageFrequency(\levelParameters{\level})
	\cdot
	\prod_{1\leq i \leq \level} \indicatorm^i_{\parameterAccess{\level}{\minimizerLength}{i}} \indicatort^i_{\parameterAccess{\level}{\threshold}{i}}
	& \mid \forall l \in \setOfLevels
\end{align*}

\subsection{Lookup Time}
The search time on a level $l \in \setOfLevels$ can be modeled analogously to that on one level:
\begin{align*}
	\timeForOperationOnStructure{level}{}^\level & \in \RR^+_0 &\mid \forall l \in \setOfLevels\\
	\timeForOperationOnStructure{level}{}^\level &= 2\timeForOperationOnStructure{\select}{S} + \frac{\varAverageFrequency^\level+1}{2}  \cdot (\timeForOperationOnStructure{\access}{O} + \timeForOperationOnStructure{\compare}{T}) &\mid \forall l \in \setOfLevels
\end{align*}

For $\varCDF{K}^\level$ of the k-mers on a level $l \in \setOfLevels$ the lookup is completed by a search through the offsets on that level whilst the remaining $(1-\varCDF{K}^\level)$ of queries result in a lookup on the next level.\\
We can thus model the lookup time recursively:
\begin{align*}
	\timeForOperationOnStructure{lookup}{}^{l} &= 
	\begin{cases}
	\begin{aligned}
		\timeForOperationOnStructure{\contains}{R} &+ \varCDF{K}^\level \cdot \timeForOperationOnStructure{level}{}^\level\\
		&+ (1-\varCDF{K}^\level) \cdot \timeForOperationOnStructure{lookup}{}^{l+1}
	\end{aligned} &\mid l \in \setOfLevels\\
	\timeForOperationOnStructure{check}{H} &\mid l \notin \setOfLevels
\end{cases}
\end{align*}

Considering the set $\setOfLevels = \{ 1,2 \}$ of two levels as an example, this results in the following total lookup time:
\begin{align*}
	\timeForOperationOnStructure{lookup}{}^1 = \timeForOperationOnStructure{\contains}{R} +&~ \varCDF{K}^1 \cdot \timeForOperationOnStructure{level}{}^1\\
		+&~ (1-\varCDF{K}^1) \cdot \Big( \timeForOperationOnStructure{\contains}{R} + \varCDF{K}^2 \cdot \timeForOperationOnStructure{level}{}^2\\
		&~ \phantom{(1-\varCDF{K}^1) \cdot \Big( \timeForOperationOnStructure{\contains}{R^2}} + (1-\varCDF{K}^2) \cdot \timeForOperationOnStructure{check}{H} \Big)\\
\end{align*}
Inserting $\timeForOperationOnStructure{level}{}^1$ and $\timeForOperationOnStructure{level}{}^2$ yields:
\begin{align*}
	\timeForOperationOnStructure{lookup}{}^1 =
		\timeForOperationOnStructure{\contains}{R} +&~ \varCDF{K}^1 \cdot (2\timeForOperationOnStructure{\select}{S} + \frac{\varAverageFrequency^1+1}{2}  \cdot (\timeForOperationOnStructure{\access}{O} + \timeForOperationOnStructure{\compare}{T}))\\
			+&~ (1-\varCDF{K}^1) \cdot \Bigg(\timeForOperationOnStructure{\contains}{R} + \varCDF{K}^2 \cdot (2\timeForOperationOnStructure{\select}{S} + \frac{\varAverageFrequency^2+1}{2}  \cdot (\timeForOperationOnStructure{\access}{O} + \timeForOperationOnStructure{\compare}{T}))\\
			&~\phantom{(1-\varCDF{K}^1) \cdot \Bigg(\timeForOperationOnStructure{\contains}{R}} + (1-\varCDF{K}^2) \cdot \timeForOperationOnStructure{check}{H} \Bigg)
\end{align*}
\todo[inline]{Frage: Soll ich das Einsetzten lieber weglassen?}

\subsection{Space Consumption}
\textbf{The text $T$} is only being stored once, so the constant approximation $bpk_T \approx 2$ can be used for several levels as well.\\

\textbf{The number of distinct \infrequent minimizers} $\varNumberOfDistinctMinimizers^\level$ on each level $\level$ can modeled analogously to $\varCDF{M}^\level$, $\varCDF{K}^\level$ and $\varAverageFrequency^\level$:
\begin{align*}
	\varNumberOfDistinctMinimizers &\in \RR^+_0 &\mid \forall l \in \setOfLevels
	\\
	\varNumberOfDistinctMinimizers &= \sum\limits_{\levelParameters{\level} \in \possibleMinimizerLengths^\level \times \possibleThresholds^\level}
	\numberOfDistinctMinimizers{\levelParameters{\level}} 
	\cdot
	\prod_{1\leq i \leq \level} \indicatorm^i_{\parameterAccess{\level}{\minimizerLength}{i}} \indicatort^i_{\parameterAccess{\level}{\threshold}{i}}
	& \mid \forall l \in \setOfLevels
\end{align*}

\textbf{Vectors $R^\level$, $S^\level$ and $O^\level$:} The bits per k-mer needed to store the bit vectors $R^\level$ and $S^\level$ as well as the offset vector $O^\level$ on every level $\level$ can be modeled analogously to using only one level:
\begin{align*}
	bpk_{R^\level} & \in \RR^+_0 &\mid \forall l \in \setOfLevels
	\\
	bpk_{S^\level} & \in \RR^+_0 &\mid \forall l \in \setOfLevels
	\\
	bpk_{O^\level} & \in \RR^+_0 &\mid \forall l \in \setOfLevels
	\\
	bpk_{R}^\level &= \frac{\varNumberOfDistinctMinimizers^\level}{\nkmers} \cdot \left(\lceil 2\varMinimizerLength^\level - log_2\varNumberOfDistinctMinimizers^\level \rceil + 3 \right) &\mid \forall l \in \setOfLevels
	\\
	bpk_{S^\level} &= \frac{\varNumberOfDistinctMinimizers^\level \cdot \varAverageFrequency^\level}{\nkmers} &\mid \forall l \in \setOfLevels
	\\
	bpk_{O^\level} &= bpk_{S^\level} \cdot log_2 |T| &\mid \forall l \in \setOfLevels
\end{align*}

\textbf{The hash set $H$} only contains an entry for k-mers $\kmer$, which have frequent minimizers on all levels.
The bits per k-mer for the hash set can thus be modeled as 
\[ bpk_H = \left(\prod_{l=1}^{L} 1-\varCDF{K}^\level\right) \cdot \frac{64}{\alpha}\]\\

\textbf{The total bits per k-mer} for $L$ levels only differ to those on one level, as there are now a vector $R,S$ and $O$ on each level:
\[ bpk_{total} = bpk_T + \sum_{l=1}^{L} \left(bpk_{R^\level} + bpk_{S^\level} + bpk_{O^\level}\right) + bpk_H\]

\chapter{Experiments}\label{chap:experiments}
\todo[inline]{aktualisieren}
We evaluate how suited the MINLP proposed in \cref{chap:minlp} is for finding optimal level parameters by analyzing how accurate the modeled lookup time and space consumption are compared to the actual lookup time measured and space needed when building RSHash. We then examine how well the model can be optimized using the SCIP solver \cite{AchterbergBertholdKochetal.2008, Achterberg2009, BestuzhevaChmielaMuellerSerranoVigerskeWegscheider2023}. Whether the resulting parameters are actually optimal in practice is highly dependent on how accurate the model is in the first place.

\section{Datasets}
We use the same datasets as \citet{pibiri_sparse_2022}. Those are the whole genomes of the atlantic cod (\textit{Gadus morhua}), the common kestrel (\textit{Falco tinnunculus}) and the whole GRCh38 human genome (\textit{Homo sapiens}) This choice allows a direct comparison with \citet{schulte-mattler_rshash_2026} who use the same datasets for benchmarking RSHash.\\
We assess the proposed MINLP both on the datasets themselves as well as on preprocessed versions without duplicate k-mers. For this a compacted de Bruijn Graph is built with BCALM2 \cite{chikhi_compacting_2016} whose maximal unitigs (non-branching paths) are extracted using the tool UST \cite{doi:10.1089/cmb.2020.0431}.\\


\section{Considered Parameters} \label{sec:consideredParameters}
For the following experiments we use k-mers of length $k = 31$ and the largest number of levels $\numberOfLevels = 3$ that is currently supported by RSHash.\\
RSHash currently works with thresholds up to $255$.
Technically, using the proposed MINLP to find optimal level parameters for a specified optimization goal, one would have to consider every minimizer lengths $1 \leq \minimizerLength < k$ and all thresholds $1 \leq \threshold \leq 255$. However, using the proposed model requires computing the distributions of minimizers for each possible combination of parameters. In practice using $3$ levels would require computing the distributions for $30^3 \cdot 255^3 \approx 448 \cdot 10^9$ combinations of level parameters which is is not viable.\\

When using the maximum threshold $\threshold=255$ and minimizers of length $15$ leads to a percentage of \infrequent  minimizers below $percentage$\todo{percentage einfügen!!} for each of the datasets considered. As this already leads to a comparatively large hash set, minimizer lengths $\minimizerLength < 15$ can be disregarded when searching for level parameters which lead to relatively low query times whilst also keeping the data structure fairly compact.\\
When using a small threshold $\threshold=32$ the percentage of \infrequent minimizers of length $25$ is above $percentage$\todo{percentage einfügen!!} for each of the datasets considered. 
Whilst this would lead to a small hash set, the compression would be small, as each minimizer covers only approximately $(31-25+2)/2 = 4$ k-mers\todo{begründung überarbeiten; approximation gilt nur für erstes level und eventuell gar nicht für unitigs}. Because of this, considering minimizer lengths between fifteen and twenty five should be sufficient for analyzing the model. To further reduce the number of parameter combinations we only consider uneven minimizer lengths.
\begin{equation}
	\possibleMinimizerLengths = \{\minimizerLength \in \NN \mid 15 \leq \minimizerLength \leq 25 \wedge \minimizerLength \% 2 = 1\}
\end{equation}
Recall that cumulative distribution functions $\CDF{\compactZV{M}{\levelParameters{\level}}}$ and $\CDF{\compactZV{K}{\levelParameters{\level}}}$ have logarithmic growth\todo{zeigen/beschreiben!!}. It follows that using exponentially increasing thresholds should be sufficient for analyzing the MINLP.
We consider the thresholds $32,64$ and $128$.\todo{reason}

\begin{equation}
	\possibleThresholds = \{32,64,128\}
\end{equation}
With these sets of considered parameters, using $3$ levels would result in $6^3 \cdot 3^3 = 5,832$ combinations of level parameters. We reduce this even further by considering only non-decreasing parameter values on subsequent levels for both the minimizer lengths and thresholds. That means each combination of level parameters $\levelParameters{3} = (\minimizerLength^1, \minimizerLength^2, \minimizerLength^3, \threshold^1, \threshold^2, \threshold^3)$ fulfills $\minimizerLength^i \leq \minimizerLength^i+1$ and $\threshold^i \leq \threshold^i+1$ for all $i \in \{1,2\}$ which can be done because\todo{reason}. \todo[inline]{explain the constraint enforcing this in the MINLP and how the sum over possible minimizer combinations changes}
\begin{align}
	constraints...
\end{align}
For a non-decreasing sequence of length $k'$, chosen from $n$ values, the number of possibilities is $\binom{n+k'-1}{k'}$. That means that when choosing three minimizer lengths out of six possible values and three thresholds out of three possible values, the total number of possible parameter combinations is $\binom{6+3-1}{3} \cdot \binom{3+3-1}{3} = 560$.

\section{Computing the Model Input} \label{sec:computationOfModelInput}
In order to evaluate the accuracy of the proposed MINLP, we first show how to compute the model input for each of the raw/preprocessed datasets which we refer to as the set of sequences $\sequences{0}$. This is done for $\numberOfLevels = 3$ levels and non-decreasing level parameters as described in \cref{sec:consideredParameters}.\\
The required stats are the total number $\nkmers$ of k-mers in the input sequence and the cumulative distributions of minimizer frequencies $\CDF{\compactZV{M}{\levelParameters{\level}}}$ and $\CDF{\compactZV{K}{\levelParameters{\level}}}$, the average frequencies $\averageFrequency(\levelParameters{\level})$ and the number of distinct \infrequent minimizers $\numberOfDistinctMinimizers{\levelParameters{\level}}$ for every non-decreasing combination of level parameters $\levelParameters{\level} \in \possibleMinimizerLengths^\level \times \possibleThresholds^\level$ on every level $1 \leq \level \leq 3$.\\

These stats can be brute-forced for every dataset considered by iterating over each non-decreasing combination of the specified level parameters and actually building RSHash with those parameters. We use the RSHash algorithm implemented by \cite{schulte-mattler_rshash_2026} which is available at \url{https://github.com/jonsmcode/rshash}.\\ Whilst there are more efficient approaches to only compute the model input, we chose this approach as it not only computes an accurate model input but also allows to measure space consumption of the data structures components and the operation times on each of these components which we will later use to assess the accuracy of the values modeled by the MINLP.\\
Specifically we measure the bits per k-mer of the vectors $R, S$ and $O$ on each level as well as the bits per k-mer to store the text and hash set which are only stored once and not per level. We also measure the average lookup time $\timeForOperationOnStructure{lookup}{}$ by measuring the lookup for $100000$ k-mers randomly selected from the sequence.\\

\begin{algorithm} 
	\caption{Computation of Minimizer Distribution} \label{algo:computeMinDistr}
	\begin{algorithmic} 
		\Input $
		\text{sequences } \sequences{0},
		\text{max Level } \numberOfLevels,
		\text{considered minimizer lengths and thresholds } \possibleMinimizerLengths, \possibleThresholds
		$
		\Function{statsComputation}{$\sequences{0}, \numberOfLevels, \possibleMinimizerLengths, \possibleThresholds$}
		\State $\nkmers \gets ccountKmers(\sequences{0})$
		\For{$\levelParameters{L} \in \possibleMinimizerLengths^L \times \possibleThresholds^L$}
		\If{$\levelParameters{L}.isDecreasing()$} \State \textbf{continue} \EndIf
		\State $RSHash \gets buildWithStats(\sequences{0}, \levelParameters{L})$
		\State $\CDF{\compactZV{M}{\levelParameters{\level}}},
		\CDF{\compactZV{K}{\levelParameters{\level}}},
		\averageFrequency(\levelParameters{\level}), \numberOfDistinctMinimizers{\levelParameters{\level}} \gets RSHash.getDistributionStats()$
		\State $spaceConsumptions [\levelParameters{L}] \gets RSHash.getSpaceConsumptions()$
		\State $\timeForOperationOnStructure{lookup}{} [\levelParameters{L}] \gets RSHash.benchLookup(noSamples=100000)$
		\EndFor
		\EndFunction
	\end{algorithmic}
\end{algorithm}
Lastly, we use a constant approximation for the times needed for a contains query on bitvector $R$, a select queries on bitvector $S$, accessing an offset in $O$ a comparison with the text $T$ at a random offset and checking whether a key exists in the hash set $H$. Based on measurements we use:
\begin{align}
	\timeForOperationOnStructure{\contains}{R^l} &= 30ns\\
	\timeForOperationOnStructure{\select}{S^l} &= 35ns  \\
	\timeForOperationOnStructure{\access}{O^l} &= 100ns  \\
	\timeForOperationOnStructure{\compare}{T} &= 100ns   \\
	\timeForOperationOnStructure{check}{H} &= 20ns
\end{align}
\todo[inline]{welche werden final genuzt?}


\section{Observations Regarding the Computed Stats}
\subsection*{Minimizer Distribution}
\begin{figure} \label{fig:multilevelDistribution}
	\centering
	\begin{tikzpicture}
	\begin{groupplot}[
		group style={
			group size=3 by 1,
			horizontal sep=0.07\textwidth,
		},
		width=0.33\textwidth,
		height=5cm,
		ymin=0, ymax=1,
		xlabel={$f$},
		]
		
		% First panel
		\nextgroupplot[
			title={Kestrel},
			ylabel={$F_M$},
			legend to name=sharedlegend,
			legend columns=3,
			]
		\addplot table[col sep=comma, x=frequency, y=F_M_1]
		{./data/F_M_Falco_tinnunculus.FalTin1.0.dna.toplevel.csv};
		\addlegendentry{Level $1$}
		\addplot table[col sep=comma, x=frequency, y=F_M_2]
		{./data/F_M_Falco_tinnunculus.FalTin1.0.dna.toplevel.csv};
		\addlegendentry{Level $2$}
		\addplot table[col sep=comma, x=frequency, y=F_M_3]
		{./data/F_M_Falco_tinnunculus.FalTin1.0.dna.toplevel.csv};
		\addlegendentry{Level $3$}
		
		% Second panel
		\nextgroupplot[title={Cod}]
		\addplot table[col sep=comma, x=frequency, y=F_M_1]
		{./data/F_M_Gadus_morhua.gadMor3.0.dna.toplevel.csv};
		\addplot table[col sep=comma, x=frequency, y=F_M_2]
		{./data/F_M_Gadus_morhua.gadMor3.0.dna.toplevel.csv};
		\addplot table[col sep=comma, x=frequency, y=F_M_3]
		{./data/F_M_Gadus_morhua.gadMor3.0.dna.toplevel.csv};
		
		% Third panel
		\nextgroupplot[title={Human}]
		\addplot table[col sep=comma, x=frequency, y=F_M_1]
		{./data/F_M_Homo_sapiens.GRCh38.dna.toplevel.csv};
		\addplot table[col sep=comma, x=frequency, y=F_M_2]
		{./data/F_M_Homo_sapiens.GRCh38.dna.toplevel.csv};
		\addplot table[col sep=comma, x=frequency, y=F_M_3]
		{./data/F_M_Homo_sapiens.GRCh38.dna.toplevel.csv};
		
	\end{groupplot}
	
	\node at (group c2r1.south) [below=1.2cm] {\pgfplotslegendfromname{sharedlegend}};
\end{tikzpicture}

	\caption{The cumulative percentage $F_M$ of minimizers with frequency up to $f$ on each level for the raw genomes Kestrel, Cod and Human. The values were measured using minimizers of length $17$ and the threshold $32$ for every level.}
\end{figure}
Recall that RSHash uses several levels aiming to repeatedly take advantage of the sparse and skewed distribution of minimizers. This requires using a different hash functions to compute minimizers on each level as using the same function on every level would mean that give equal minimizer length and threshold, those minimizers which are \frequent on a level would also be \frequent on the next.
We observe that despite of the minimizers being computed using a different hash function, their distributions on levels $\level > 1$ remain highly dependent on the distributions on all previous levels.  \Cref{fig:multilevelDistribution} shows this for the raw genomes using the level parameters $\levelParameters{\level} = (17,17,17,32,32,32)$.\\

It seems that super-k-mers with \frequent minimizers on one level are more likely to also have \frequent minimizers on subsequent levels.
A possible explanation for this observation are repetitively regions of the sequence in which a subsequences shorter than $k$ repeats many times over. K-mers in such regions would be likely to share the same minimizer regardless of the hash function used to compute it, resulting in those minimizers being frequent on every level. This could be verified by searching for such repeating substrings in those super-k-mers which are frequent on all levels.

\subsection*{Operation Times}
Recall that times 
$\timeForOperationOnStructure{\contains}{R}$,
$\timeForOperationOnStructure{\select}{S}$,
$\timeForOperationOnStructure{\access}{O}$,
$\timeForOperationOnStructure{\compare}{T}$,
$\timeForOperationOnStructure{check}{H}$
needed for a contains rank query on bitvector $R$, a select queries on bitvector $S$, accessing an offset in $O$, a comparison with the text $T$ at a random offset and checking whether a key exists in the hash set $H$ are part of the model input and considered to be constant.\\
The measured times for these operations show that they actually vary in dependence of the level parameters. This is expected as ...

\section{Model Accuracy}
We asses the accuracy of the lookup time and space consumption modeled by the MINLP for each of the datasets considered. To do this we first compute model input, the actual space consumption and the measured operation times for every non-decreasing combination of level parameters on three levels as outlined in \cref{sec:computationOfModelInput}.\\
For each raw and preprocessed sequence and each combination of the level parameters considered, we then use a modified version of the proposed model which forces the exact combination of level parameters and returns the modeled space consumption and lookup time.

\subsection{Lookup Time}


\subsection{Space Consumption}
\begin{tabular}{l|cccccccccccc}
	\bfseries metric & \bfseries bpk\_T & \bfseries bpk\_R\_0 & \bfseries bpk\_R\_1 &
	\bfseries bpk\_R\_2 & \bfseries bpk\_S\_0 & \bfseries bpk\_S\_1 & \bfseries bpk\_S\_2 &
	\bfseries bpk\_O\_0 & \bfseries bpk\_O\_1 & \bfseries bpk\_O\_2 & \bfseries bpk\_H &
	\bfseries bpk\_total \\
	\hline
	\csvreader[head to column names]{./data/stats_evaluating_space_consumptionFalco_tinnunculus.FalTin1.0.dna.toplevel.csv}{}{%
		\csname metric\endcsname & \csname bpk\_T\endcsname & \csname bpk\_R\_0\endcsname &
		\csname bpk\_R\_1\endcsname & \csname bpk\_R\_2\endcsname &
		\csname bpk\_S\_0\endcsname & \csname bpk\_S\_1\endcsname &
		\csname bpk\_S\_2\endcsname & \csname bpk\_O\_0\endcsname &
		\csname bpk\_O\_1\endcsname & \csname bpk\_O\_2\endcsname &
		\csname bpk\_H\endcsname & \csname bpk\_total\endcsname\\\hline
	}
\end{tabular}

\todo[inline]{Absatz zu $bpk_{total}$ schreiben}
For both the raw and preprocessed datasets the modeled space consumption of the vectors $R^\level, S^\level$ and $O^\level$ on every level $1 \leq \level \leq 3$ is accurate up to an error of $..$\todo{max error eintragen}.\\
The only significant deviations are for the bits per k-mer $bpk_T$ and $bpk_H$ needed to store the text and the hash set $H$. The chosen approximation of $bpk_T \approx 2$ holds for the raw sequences because those sequences are generally much longer than $k$. This is not necessarily the case when the input sequences have been preprocessed to remove duplicate k-mers. The preprocessed sequences are represented through unitigs, which are non-branching paths in the compact de Bruijn graph of the sequence. For highly repetitive genomes, their lengths are much closer to $k$, requiring more bits per k-mer to store. This effect is negligent for the common kestrels genome which contains almost no duplicate k-mers. The preprocessed genomes of atlantic cods and human need approximately $2.287$ and $2.312$ bits per k-mer. This could be accounted for by including the bits per k-mer needed to store a sequence in the model input as its computation is simple compared to the other stats which have to be computed for the model input.\\\todo[inline]{Frage: soll ich den Abschnitt bezüglich $bpk_T$ hier in der auswertung lassen oder die $bpk_T$ als Parameter zum Model hinzufügen (diese zu berechnen ist im Vergleich mit den restlichen stats sehr günstig)?}

RSHash is implemented in c++ using $gtl::flat\_hash\_set$ which uses one byte of metadata for each entry in the hash set to allow faster queries. For a direct comparison of the measured and modeled space consumption, we can account for this by modifying the constraint to be $bpk_H = \left(\prod_{l=1}^{3} 1-\varCDF{K}^\level\right) \cdot \frac{64 + 8}{\alpha}$ which is equivalent to adding a bias term of $(1-\CDF{\compactZV{K}{\levelParameters{1}}})(1-\CDF{\compactZV{K}{\levelParameters{2}}})(1-\CDF{\compactZV{K}{\levelParameters{3}}}) \cdot 8/0.75$ to the modeled values.\\
However, even with this adjustment the modeled space consumption of the hash set is neither correct for the raw genomes nor for the preprocessed ones. Whilst we model the bits per k-mer $bpk_H$ of the hash set as a continues function, our measurements seem to take values from a discrete distribution.\\
This observation can also be explained by the hash sets implementation: Upon declaration a c++ hash set starts of with a hash function mapping keys to a hash space of a relatively small size. That size is doubled whenever a certain criteria regarding the hash sets load is no longer fulfilled. This discrete behavior with varying loads cannot be captured by a continues model using a constant load factor.\\
In addition to this, the model tends to overestimate the hash sets size for the raw sequences. This is likely because repeating k-mers can count towards the percentage of k-mers with \infrequent minimizers on all levels and thus lead to a larger value of the modeled $bpk_H$. However, in practice these duplicate k-mers collapse to single entry in the hash set as the lookup only requires determining a single occurrence. This effect is again negligent for the common kestrels genome, where
\todo[inline]{Plots? Wenn ja, welche?}

\section{Limitations}
The poorly modeled lookup time renders the MINLP useless for finding optimal level parameters for a specified optimization goal.

This approaches plausibility is heavily dependent on the time needed to compute the model input as outlined \cref{sec:computationOfModelInput}. Especially the runtime for the computation of the distributions $\probabilityDistribution{\ZVM{\level}}$ and $\probabilityDistribution{\ZVK{\level}}$ for every combination of the level parameters $\levelParameters{\level}$ on every level $\level$ grows exponentially with the number of levels when computed by building RSHash.

For this reason, an approach yielding only an approximation of the optimal level parameters but taking less time to do so could be of interest.

We will further explore methods \cref{chap:subsampling} the genome in order to reduce the computation of minimizer distribution.\todo{TODO}

In another approach we will try to approximate the distributions on higher levels based on that on the first level in order to get around the exponential cost of computing those values.\todo{TODO}

\chapter{Outlook} \label{chap:outlook}
\todo[inline]{model negative lookups with a parameter for the ratio of positive/negative lookups}

\textbf{Locate}
\todo[inline]{outline differences to lookup}
\begin{itemize}
	\item lookup time: check every offset instead of only an average of half of the offsets
	\item hash table instead of hash set - contains one entry (key:k-mer, value:offset) for each appearance of every k-mer (no longer distinct) - for the unitigs, this shouldn't have an effect
	\item Locate benötigt eine Liste mit allen offsets: $bpk_H = (1-\varCDF{K}) \cdot \frac{(64 + (log_2|T|\cdot frequency))}{\alpha}$
\end{itemize}

\textbf{Extrapolation Minimizer Distributions:} As the minimizer distribution on each level is affected by that on the previous ones, it might be possible to only compute a selection of distributions and extrapolate the remaining ones. This 

\textbf{Sub sampling}
For large genomes, the computing of the model input could be sped up  by randomly sub sampling a certain percentage of the genomes k-mers and computing the minimizer distributions on those. A smaller percentage would then mean a faster computation at the cost of losing accuracy.

\chapter{Conclusion}

\chapter{stuff}
A fairly straightforward way\todo{zu umgangssprachlich?} of computing these is to use a recursion over the levels. The input for each recursive call consist of the current level $\level$, the number of levels $\numberOfLevels$, the previous level parameters $\levelParameters{\level-1}$ and the set of input sequences in dependence of those parameters.\\
Upon being called, the first minimizer $\minimizerLength \in \possibleMinimizerLengths$ is selected and used to compute the distribution stats for the input sequence on this level in dependence of $\minimizerLength$. The function then loops over each threshold $\threshold \in \possibleThresholds$, removes all k-mers with \infrequent minimizers when using $(\minimizerLength, \threshold)$ from the input sequences and then calls the recursion for the next level, which receives the next level, the previous parameters extended by $(\minimizerLength, \threshold)$ and the modified sequences as input.
Only once this recursion is completed for every threshold $\threshold \in \possibleThresholds$, the next minimizer is picked from $\possibleMinimizerLengths$. The described algorithm:
\todo[inline]{soll ich hier erklären, wie die Berechnung in der Praxis noch ein bisschen schneller geschieht? Oder ist das an dieser Stelle egal, wenn es am grundlegenden Problem der exponentiellen Laufzeit nichts ändert?}
\begin{algorithm} 
	\caption{Computation of Minimizer Distribution}
	\begin{algorithmic} 
		\Input $
		\text{level } \level,
		\text{max Level } \numberOfLevels,
		\text{level parameters } \levelParameters{\level-1},
		\text{sequences } \sequences{\level-1}
		$
		\Function{recursiveMinimizerDistribution}{$\level, \numberOfLevels, \levelParameters{\level-1} \sequences{\level-1}$}
		\If{$\level = \numberOfLevels$} \State \textbf{return} \EndIf
		\For{$\minimizerLength \in \possibleMinimizerLengths$}
		\For{$\threshold \in \possibleThresholds$}
		\State $\levelParameters{\level} \gets extend(\levelParameters{\level-1}, m,t)$	\Comment $extend((m^1,t^1), m^2,t^2) = (m^1,m^2,t^1,t^2)$
		\State $\sequences{\level} \gets getFrequentKmers(\level, \sequences{\level-1}, \minimizerLength, \threshold)$
		\State 
		\State $\CDF{\compactZV{M}{\levelParameters{\level}}},
		\CDF{\compactZV{K}{\levelParameters{\level}}},
		\averageFrequency(\levelParameters{\level}), \numberOfDistinctMinimizers{\levelParameters{\level}} \gets distributionStats(\hashfunction{\level}, \sequences{\level-1}, m)$
		\State 
		\Call{recursiveMinimizerDistribution}{$\level+1, \numberOfLevels, \levelParameters{\level}, \sequences{\level}$}
		
		\EndFor
		\EndFor
		\EndFunction
	\end{algorithmic}
\end{algorithm}

This approach is guaranteed to compute correct minimizer distributions for every set of level parameters on every level as it effectively performs the same steps as the RSHash algorithm only without building the data structure itself. This also means, that one could include the building of bitvectors $R,S$ and $O$ on every level as well as the hash set $H$ and simply measure the space consumption and lookup time for every combination, which would render the MINLP obsolete.\\
However, this approach is sufficient for computing the input needed for analyzing how well the MINLP models the lookup time and space consumption using only the minimizer distributions. Possibilities of speeding up the time needed are outlined in \cref{chap:outlook}.

\printbibliography[heading=bibintoc,title={Bibliography}]

\end{document}
